\section{Consideraciones éticas}
\label{sec:consideraciones_eticas}

La investigación involucra como participantes a estudiantes menores de edad (15 a 17 años), lo que exige el cumplimiento riguroso de procedimientos éticos orientados a proteger su integridad, su privacidad y su participación libre e informada. Los procedimientos descritos a continuación se aplican durante todas las fases del estudio.

\subsection{Consentimiento y asentimiento informado}

Antes del inicio de cualquier actividad de recolección de datos se gestionan tres documentos de consentimiento:

\begin{itemize}
\item \textbf{Consentimiento de padres o acudientes:} documento escrito dirigido al acudiente de cada estudiante, que explica los objetivos del estudio, las actividades en las que participará el estudiante, el tipo de datos que se recolectarán (observaciones, producciones escritas, grabaciones de audio o video, entrevistas), y el uso exclusivamente investigativo de esa información. La firma del documento es condición necesaria para que el estudiante participe en las actividades registradas.

\item \textbf{Asentimiento del estudiante:} documento escrito dirigido al propio estudiante, redactado en lenguaje accesible, que le explica en qué consiste su participación, que puede retirarse en cualquier momento sin consecuencias para su desempeño académico, y que tiene derecho a no responder preguntas que no desee responder en las entrevistas.

\item \textbf{Autorización institucional:} comunicación formal dirigida al rector o coordinador académico del Colegio Villa de las Palmas, que informa sobre el propósito de la investigación, el período de implementación, y el compromiso de compartir los resultados con la institución al finalizar el estudio.
\end{itemize}

\subsection{Confidencialidad y anonimato}

La identidad de los estudiantes no aparece en ningún reporte, análisis o publicación derivada de esta investigación. Para proteger el anonimato:

\begin{itemize}
\item Cada estudiante es identificado mediante un código alfanumérico (E01, E02\ldots E25) que se asigna al inicio de la implementación y se usa de forma consistente en todos los documentos de análisis.
\item Los registros de audio, video y producciones escritas se almacenan en dispositivos de acceso restringido, protegidos con contraseña, a los que solo tiene acceso el investigador principal.
\item Los nombres de los estudiantes no aparecen en transcripciones, notas de campo ni materiales de análisis.
\end{itemize}

\subsection{Manejo de grabaciones}

Las sesiones de trabajo podrán ser grabadas en audio o video únicamente cuando exista consentimiento firmado del acudiente y asentimiento del estudiante. Las grabaciones:

\begin{itemize}
\item Se utilizan exclusivamente para los fines de esta investigación (transcripción, análisis de procesos de razonamiento).
\item No se publican ni comparten en plataformas digitales.
\item Serán eliminadas de todos los dispositivos de almacenamiento en un plazo no mayor a dos años después de la defensa del trabajo de grado.
\end{itemize}

\subsection{Participación voluntaria y sin consecuencias académicas}

La participación en el estudio es voluntaria. Un estudiante puede:

\begin{itemize}
\item Decidir no ser grabado sin que ello afecte su participación en las actividades de clase.
\item Retirarse de las entrevistas o de cualquier actividad de recolección de datos en el momento que lo decida.
\item Solicitar que sus producciones escritas no sean incluidas en el análisis.
\end{itemize}

En ningún caso la decisión de no participar en actividades de recolección de datos afectará la calificación del estudiante en la asignatura de matemáticas.

\subsection{Riesgo mínimo}

Esta investigación se clasifica como de \textbf{riesgo mínimo}, de acuerdo con los criterios establecidos para investigación con seres humanos en Colombia. Las actividades que se desarrollan (resolución de tareas matemáticas, uso de software educativo, conversación con un sistema de IA en contexto académico) no representan riesgo para la integridad física, psicológica o social de los participantes. No se realiza ninguna intervención clínica, diagnóstica ni invasiva. El investigador toma las precauciones necesarias para que el ambiente de trabajo sea seguro, respetuoso y libre de presión.
